\section{Large Language Models w inżynierii oprogramowania - kontekst technologiczny}
\label{subsec:llm-kontekst-technologiczny}

Duże modele językowe (Large Language Models, LLMs) stanowią aktualnie jedną z najważniejszych klas modeli w kontekście
informatyki wykorzystywanych w zadaniach związanych z wytwarzaniem oprogramowania. Ich istotna rola wynika przede wszystkim
z faktu, że kod źródłowy, tak jak język naturalny, ma strukturę sekwencyjną oraz jest silnie determinowany przez lokalny i globalny
kontekst (zależność między innymi plikami, konwencje nazewnictwa, wzorce projektowe i architektoniczne).
Sprawia to, że modele uczone na dużych blokach tekstu i kodu mogą sprawnie wspierać zadania wymagające
jednocześnie zrozumienia kontekstu i intencji programisty, jak również generowania precyzyjnych komentarzy czy sugestii zmian.

W kontekście tej pracy magisterskiej istotne jest zrozumienie, jak mechanizmy technologii informatycznej umożliwiają użycie LLM w potokach
pracy programisty oraz z jakimi ograniczeniami należy się liczyć. Podczas projektowania systemu automatycznego
przeglądu kodu kluczowe znaczenie mają: paradygmat prompt engineering jako warstwa sterowania zachowaniem modelu,
techniki wzbogacania kontekstu (w szczególności RAG) oraz metody zwiększania stabilności i powtarzalności wyników.

\subsubsection{Prompt engineering jako warstwa sterowania zachowaniem LLM}
W zastosowaniach związanych z wytwarzaniem oprogramowania duże modele językowe rzadko pracują w trybie "surowy", czyli
bez jawnie zdefiniowanych instrukcji. Najczęściej spotykaną praktyką jest formułowanie zadań w postaci promptu,
który określa jednocześnie cel (np. wskazanie problemów w kodzie), jak i oczekiwaną  formę odpowiedzi (np. lista uwag z
priorytetem, uzasadnieniem i propozycją poprawki). Systematyczny przegląd literatury Syahputri i in. dotyczący prompt
engineering w obszarze inżynierii oprogramowania wskazuje, że badania w latach 2020-2025 koncentrowały się na czterech
dominujących obszarach: ręcznym tworzeniu promptów, podejściu retrieval-augmented generation, chain-of-thought prompting
oraz automatycznym strojeniu promptów \cite{syahputri2025prompt}. Autorzy podkreślają w publikacji także typowe problemy
praktyczne, takie jak \emph{prompt brittleness} (wrażliwość na drobne zmiany treści polecenia), brak standardowych
protokołów ewaluacji oraz wyzwania skalowalności \cite{syahputri2025prompt}.

Z punktu widzenia projektanta systemu automatycznego przeglądu kodu, prompt engineering docelowo pełni rolę porównywalną
do konfiguracji reguł w klasycznych narzędziach statycznej analizy kodu. Określa co będzie sprawdzane (np. bezpieczeństwo
wątków, czytelność, zgodność ze standardami projektu) oraz jak wynik zostanie zakomunikowany (np. komentarz). W odróżnieniu
od standardowych podejść regułowych, odpowiednie prompty mogą rzutować na bardziej złożone oczekiwania wobec semantyki i
intencji zmian, np. pomagać weryfikować zgodność z architekturą warstwową lub spójność logiki biznesowej. Oznacza to, że
jakość automatycznej recenzji w dużej mierze zależy od jakości i powtarzalności promptów, a nie wyłącznie od możliwości
użytego modelu.

\subsubsection{Stabilność wyników: few-shot, prompt chaining i metody redukcji niedeterministyczności}
Analizując praktyczność LLM w kontekście analizy kodu, jednym z problemów jest niedeterministyczność odpowiedzi oraz duża zależność
od sposobu sformułowania polecenia. Analiza wykorzystania ChatGPT (GPT-4-turbo) do detekcji wyścigów danych w
programach równoległych pokazuje, że proste, bazowe prompty potrafią prowadzić do nieużytecznych wyników, np. oznaczania
wszystkich plików podlegających analizie jako wadliwe \cite{alsofyani2025datarace}. Alsofyani i in. w swoim badaniu pokazali, że
techniki wprowadzające strukturę do interakcji z LLM, w szczególności few-shots oraz podejście wieloetapowe (prompt chaining)
poprawiają skuteczność klasyfikacji w porównaniu z prostymi wariantami zero-shot. Dla scenariusza z pojedynczym błędem
na plik najlepszy wariant (prompt chaining + few-shot) osiągnął \mbox{77.7\%} accuracy, a sam "basic prompt" dawał
\mbox{50.0\%} \cite{alsofyani2025datarace}. Istotnym faktem jest stabilizowanie wyników przez autorów poprzez trzykrotne
uruchamianie i zastosowanie głosowania większościowego (co najmniej 2/3 zgodnych odpowiedzi) \cite{alsofyani2025datarace}.
Adaptacja tego typu procedur jest szczególnie istotna w kontekście automatycznych komentarzy publikowanych w procesie code review.
Niepożądane wahania mogą obniżać zaufanie zespołu do narzędzia.

Podejście few-shots jest metodą przekazywania modelowi tzw. lokalnej normy jakościowej - przykładów tego, jakie uwagi
są uznawane za trafne i użyteczne w danym kontekście projektowym. Z kolei prompt chaining oznacza dekompozycję problemu
na etapy, np. (1) identyfikacja komponentów dotkniętych zmianą, (2) ocena ryzyka i zgodności ze standardami, (3) napisanie
komentarzy w oczekiwanym formacie. Takie podejście do sterowania w formie potoku jest szczególnie kompatybilne z procesem
Pull/Merge Request, gdzie wynik powinien być podawany w sposób jednoznaczny i operacyjny.

\subsubsection{Automatyzacja optymalizacji promptów}
Proces ręcznego dopracowywania promptów jest kosztowny i trudny do skalowania w dużych organizacjach (różne języki, moduły i standardy).
W odpowiedzi na ten problem pojawiły się prace traktujące prompt jako obiekt optymalizacji. W swojej pracy Ye i in.
proponują ciekawe podejście Prompt Alchemy (Prochemy), w którym prompt podlega iteracyjnemu ulepszaniu w pętli mutacja-ewaluacja-selekcja.
Ocena wariantów polega na sygnale wykonawczym (execution-driven), inaczej - przejściu testów \cite{ye2025promptalchemy}.
Podczas eksperymentów odnotowano stabilne przyrosty jakości - średnio  \mbox{+4.04\%} względem baseline zero-shot 
na benchmarkach HumanEval/MBPP oraz \mbox{+14.15\%} na LiveCodeBench \cite{ye2025promptalchemy}. W kontekście tworzenia
narzędzia do automatycznego przeglądu kodu jest to cenna obserwacja: jeśli jakość recenzji ma być mierzalna i porównywalna,
należy stosować walidacje opartą o konkretne sygnały (np. wyniki testów, linterów, analizatorów bezpieczeństwa), a nie tylko
o subiektywną ocenę jakości tekstu wygenerowanego przez LLM.

\subsubsection{Retrieval-Augmented Generation jako mechanizm uziemienia odpowiedzi w artefaktach projektu}
Z technicznego punktu widzenia faktem jest, że ograniczeniem LLM jest "wiedza" oparta o wyuczone wagi modelu.
Utrudnia to aktualizację informacji i zwiększa szansę na wygenerowanie odpowiedzi niezgodnych z konkretnym kontekstem
projektowym. Standardową odpowiedzią na ten problem jest Retrieval-Augmented Generation (RAG), czyli łączenie
generatora z mechanizmem pobierania dokumentów z zewnętrznej bazy. W ujęciu Lewisa i in. RAG stanowi hybrydę 
pamięci parametrycznej (model generatywny) oraz nieparametrycznej (indeks dokumentów), a generacja jest warunkowana 
wynikami retrievalu \cite{lewis2020rag}. Autorzy wykorzystali w eksperymentach m.in. Wikipedię z grudnia 2018 r. podzieloną
na ok 21 mln fragmentów, a model RAG (z generatorem BART-large) osiągnął wynik najlepszy do tamtej pory w zadaniach
open-domain QA (np. \mbox{44.5} EM dla Natural Questions w wariancie RAG-Sequence) \cite{lewis2020rag}.
W kontekście tematu pracy istotne jest to, że mechanizm retrieval pozwala na warunkowanie odpowiedzi aktualnym i lokalnym
kontekstem, a dokładniej dokumentacją projektu, standardami i wytycznymi kodowania, poprzednimi decyzjami architektonicznymi
oraz historią poprzednich recenzji.

W literaturze RAG częściej jest opisywany jako strategia wzmacniająca prompt engineering poprzez dołączanie kontekstu
zewnętrznego niż jako całkowicie odrębna technika \cite{syahputri2025prompt}. Takie podejście jest zgodne z praktyką
implementacyjną: w systemie code review warstwa retrieval może zostać potraktowana jako moduł przygotowujący "pakiet kontekstu" do 
promptu, a nie jako osobny model.

\subsubsection{RAG dla zadań kodowych: korzyści, koszty i ryzyko nietrafionego kontekstu}
W zastosowaniach związanych głównie z kodem, skuteczność RAG nie zależy tylko od samej idei dołączenia dokumentów, lecz
także od jakości mechanizmów retrievalu oraz sposobu integrowania kontekstu. Yang i in. analizują retrieval-augmented framework (RAF)
w kontekście generowania kodu, pokazując, że poprawa wyników jest zauważalna, lecz silnie zależna od konfiguracji \cite{yang2025ragcodegen}.
Dla datasetu CONCODE raportowane są średnie wzrosty m.in. \mbox{+6.79\%} w EM oraz 
\mbox{+8.72\%} w CodeBLEU, natomiast dla CoNaLa poprawa CodeBLEU wynosi średnio \mbox{+16.75\%} \cite{yang2025ragcodegen}.
Autorzy wskazują jednak, że prosty BM25 bywa konkurencyjny względem podejść opartych o gęste reprezentacje, a zwiększanie
liczby integrowanych fragmentów nie zawsze daje coraz większe korzyści (prawo malejących przyrostów jakości) \cite{yang2025ragcodegen}.
Oznacza to, że w narzędziu do code review, RAG powinien być projektowany jako element konfigurowalny: dobór retrievera jak i
liczba dołączanych fragmentów powinna zależeć od projektu, języka programowania i charakter zmian.

Skupiając się na cechach RAG, warto wziąć pod uwagę ograniczenia w postaci kosztów obliczeniowych oraz ryzyka "zaśmiecenia"
promptu niewłaściwymi dokumentami. Jeśli wykorzystane zostaną fragmenty niedopasowane do ocenianych zmian, model może wytwarzać
komentarze z pozoru wiarygodne, lecz nieadekwatne. Wyniki wskazujące na przypadki pogorszenia jakości przy
konkretnych konfiguracjach RAG są argumentem za wdrożeniem walidacji i monitorowania tego elementu procesu code review
(np. przez testy regresji jakości komentarzy) \cite{yang2025ragcodegen}. Dodatkowym czynnikiem jest ograniczona długość 
wejścia modelu, co wymusza selekcję, streszczanie albo priorytetyzację kontekstu.

\subsubsection{Implikacje technologiczne dla automatycznego code review}
Powyższe obserwacje prowadzą do wniosku, że wykorzystanie LLM jako element inżynierii oprogramowania ma charakter
systemowy: skuteczność nie wynika wyłącznie z charakteru dużych modeli językowych czy konkretnego modelu, lecz
z jakości warstwy sterowania (prompty), dostępności kontekstu (retrieval) oraz procedur stabilizacji i oceny wyników.
W kontekście projektu narzędzia do automatycznego przeglądu kodu istotne są trzy aspekty: (1) projektowanie promptów
jako elementów podlegających wersjonowaniu i testowaniu, (2) umiejscowanie wynikowych komentarzy w dokumentacji i 
standardach projektowych przez mechanizmy RAG, oraz (3) stosowanie metod redukujących niedeterministyczność (np. 
prompt chaining, majority voting), gdy konsekwencją błędnej odpowiedzi mogą być kosztowne decyzje 
w procesie integracji kodu \cite{alsofyani2025datarace,syahputri2025prompt}.

Cechy te sprawiają, że LLM jest naturalnym kandydatem do zastosowania w procesach code review. Ich praktyczność zależy jednak
od odpowiedniego doboru technik prompt engineering i RAG, co stanowi punkt wyjścia do przeglądu rozwiązań stricte nakierowanych 
na automatyzację recenzji.